\documentclass[11pt]{article}
\usepackage[margin=1in]{geometry}
\usepackage{amsmath,amssymb,amsthm,hyperref}
\newtheorem{theorem}{Theorem}
\newtheorem{lemma}{Lemma}
\newcommand{\R}{\mathbb R}
\newcommand{\C}{\mathbb C}
\title{Straight centroid loci for polar Poncelet polygons}
\author{Principia Math In Progress}
\date{21 September 2026}
\begin{document}
\maketitle
\begin{abstract}
For every size of the focus-circle Poncelet family inscribed in a
parabola, the vertex and area centroids of the polar polygons trace
complete lines parallel to the directrix. The proof removes every
possible pole of their horizontal coordinates on the elliptic flag
curve. In particular, finite zero-area configurations are accounted
for rather than omitted.
\end{abstract}

\section*{Task 1: the original problem}
Let
\[
 \mathcal P:\ y^2+4fx=0,\qquad
 \mathcal C:\ (x+f)^2+y^2=r^2,\qquad 0<r<f,
\]
admit a primitive Poncelet family of $N\geq3$ distinct vertices
circumscribed about the inner circle. These are the families studied
in \cite{BGR}, including all its convex cases. For a polygon $P$ in
the family, let $Q$ be the polygon whose sides are the tangents to
$\mathcal P$ at the vertices of $P$. Area and area centroid use
signed area and winding multiplicity for a self-intersecting polygon,
as in \cite{ST,CC}.

\begin{theorem}
There is a real constant $h_0$ such that
\begin{align*}
 \{C_0(Q):P\text{ finite}\}&=\{(h_0,Y):Y\in\R\},\\
 \{C_2(Q):P\text{ finite}\}&=\{(h_2,Y):Y\in\R\},
 \qquad h_2=\frac{8f^2-r^2}{12f}.
\end{align*}
Thus both centroid loci are full lines parallel to the directrix
$x=f$. Moreover, for each polygon individually,
\[
 A(Q)=-\tfrac12 A(P),\qquad
 (C_0(Q))_y=(C_0(P))_y,\qquad
 (C_2(Q))_y=(C_2(P))_y.
\]
\end{theorem}

This proves Conjecture 5 of \cite{BGR}. Its $N=3,4$ cases are already
proved there. We use the standard elliptic-curve description of the
Poncelet map from \cite{ST}, and adapt the area zero-counting idea of
\cite{CC}. The pole orders needed in this parabolic setting are
established explicitly below. We do not infer a line from a generic
conic-locus assertion.

\section*{Coordinates and polygon identities}
Write the vertices of $P$ as
\[
 P_i=(-t_i^2/(4f),t_i),\qquad i\pmod N.
\]
The tangent at parameter $t$ is $4fx+2ty-t^2=0$. Hence
\begin{equation}\label{polar}
 Q_i=\left(-\frac{t_i t_{i+1}}{4f},
                 \frac{t_i+t_{i+1}}2\right).
\end{equation}
Put
\[
 S_2=\sum_i t_i^2(t_{i+1}-t_{i-1}),\quad
 S_3=\sum_i t_i^3(t_{i+1}-t_{i-1}),\quad
 S_{32}=\sum_i t_i^3(t_{i+1}^2-t_{i-1}^2).
\]
For any cyclic list of parameters, the shoelace formulas give
\begin{align}
 A(P)&=-\frac{S_2}{8f},& A(Q)&=\frac{S_2}{16f},\label{area}\\
 (C_2(P))_y&=(C_2(Q))_y=\frac{S_3}{3S_2},&
 (C_2(Q))_x&=-\frac{S_{32}}{12fS_2},\label{moments}\\
 (C_0(Q))_y&=\frac1N\sum_i t_i,&
 (C_0(Q))_x&=-\frac1{4fN}\sum_i t_i t_{i+1}.\label{vertex}
\end{align}
These are polynomial identities for area and first moments before
division by area. For example,
\[
 \det(Q_i,Q_{i+1})
 =\frac{t_{i+1}^2(t_{i+2}-t_i)}{8f}.
\]
Summing and shifting the index proves \eqref{area}; multiplication
by $(Q_i+Q_{i+1})/6$ gives the first moments in \eqref{moments}.
All real finite polygons considered here have $A(P)\ne0$: orient
the edges with the circle on the left, and their fan triangles from
the focus have area $r\ell_i/2>0$. Consequently $A(Q)\ne0$ too.

The chord between parameters $t,s$ has equation
$4fx+(t+s)y-ts=0$. Tangency to $\mathcal C$ is equivalent to
\begin{equation}\label{chord}
 (ts+4f^2)^2=r^2\bigl(16f^2+(t+s)^2\bigr).
\end{equation}
If $u,v$ are the two neighbors of $t$, Vieta's formula therefore gives
\begin{equation}\label{vieta}
 u+v=\frac{2(r^2-4f^2)t}{t^2-r^2}.
\end{equation}
These formulas also show that the polar vertices lie on
$ (x-f)^2/r^2-y^2/(4f^2)=1$.

\section*{The flag curve and the area zeros}
Recall the standard complex flag curve
\[
 E=\{(p,\ell):p\in\mathcal P,\ p\in\ell,
                    \ \ell\text{ tangent to }\mathcal C\}.
\]
Its compactification is a smooth elliptic curve. In the coordinate
above it has a model
\[
 w^2=t^4+8f^2t^2+16f^2(f^2-r^2).
\]
The four roots on the right are distinct when $0<r<f$, so the
degree-two map to the $t$-sphere has four simple branch points.
The two points over $t=\infty$ are unramified; $t$ has a simple
pole at each. The Poncelet step $T$ is a translation of exact order
$N$ on $E$. The involution $\tau$ changing the chosen tangent
reverses a polygon and satisfies $\tau T\tau=T^{-1}$. With a fixed
point of $\tau$ as origin, these maps have the form
$T(z)=z+a$, $\tau(z)=-z$; see \cite[Section 3]{ST}.

There are exactly $2N$ flags at which one of the $t_i$ is infinite,
and at each such flag exactly one $t_i$ is infinite. To check that
these are distinct, use the real geometry: the two real components
of $E$ correspond to the two choices of direction around the nested
pair, and each projects bijectively to the real projective parabola.
The step preserves each component and has $N$ distinct vertices per
orbit. The two flags over infinity belong to different components.

Near any such flag put the unbounded parameter equal to $t$.
Its neighbors approach $r,-r$; all other parameters stay bounded.
Splitting $P$ into the triangle at that vertex and a bounded
remainder shows
\begin{equation}\label{area-infty}
 A(P)=\pm\frac{r}{4f}t^2+o(t^2).
\end{equation}
Thus $A(P)$, viewed as a meromorphic function on $E$, has exactly
$2N$ poles, each of order two, and total pole order $4N$.

Now consider the $N$ reflections $\rho_j=T^j\tau$,
$0\leq j<N$. Each is an elliptic involution $z\mapsto ja-z$,
and so has four fixed points (the four solutions of $2z=ja$).
Their fixed-point sets are disjoint: a common fixed point for two
reflections would be fixed by a nontrivial translation $T^{j-k}$.
None of these $4N$ points is an infinite-vertex flag, since its real
polygon has distinct vertices and cannot equal its reversed labeled
vertex list up to a cyclic shift when $N\geq3$.

Oriented area is invariant under $T$ and changes sign under $\tau$.
It must therefore vanish at all $4N$ reflection fixed points.
A nonconstant meromorphic function on a compact curve has the same
total zero and pole order. By \eqref{area-infty}, these are precisely
all zeros of $A(P)$, and every one is simple. Equation \eqref{area}
gives exactly the same zero divisor for $A(Q)$.

Let $M_x(Q)=\frac16\sum_i\det(Q_i,Q_{i+1})
((Q_i)_x+(Q_{i+1})_x)$ be the signed first moment. It is finite
where all $t_i$ are finite, and it too is invariant under cyclic
relabeling and changes sign under reversal. At each of the $4N$
zeros just found, $M_x(Q)$ vanishes. Since the denominator has a
simple zero, the quotient $M_x(Q)/A(Q)=(C_2(Q))_x$ has a removable
singularity there. There are no other possible finite poles.
This supplies the finite-zero check directly, following the
zero-counting method in \cite{CC} rather than assuming that a
centroid's denominator never vanishes over $\C$.

\section*{The horizontal coordinates have no poles}
First consider $(C_0(Q))_x$. In the sum in \eqref{vertex}, the only
unbounded terms near a pole parameter $t$ combine as $t(u+v)$.
Equation \eqref{vieta} makes this bounded. All other summands are
bounded. Therefore $(C_0(Q))_x$ has no poles on $E$ and is constant.

For $(C_2(Q))_x$, finite denominator zeros have already been removed.
At an infinite-vertex flag, \eqref{area-infty} and \eqref{area}
give $S_2$ a nonzero term of order $t^2$. In $S_{32}$, the summand
with unbounded central parameter is
\[
 t^3(v^2-u^2)=t^3(v-u)(v+u)=O(t^2),
\]
because $u,v$ are bounded and \eqref{vieta} gives $u+v=O(t^{-1})$.
Each of the other summands is $O(t^2)$ as well. Thus the quotient
in \eqref{moments} stays bounded. Meromorphicity makes these
singularities removable too. A holomorphic function on the compact
connected curve $E$ is constant, proving $(C_2(Q))_x=h_2$.
The two constants are real, as the family contains real polygons.

The area-centroid constant can be evaluated at infinity. Label the
neighbors so that $u\to r$, $v\to-r$. Then
\[
 \frac{S_2}{t^2}\longrightarrow-2r,\qquad
 \frac{S_{32}}{t^2}\longrightarrow
 -4r(r^2-4f^2)+r^3-(-r)^3=2r(8f^2-r^2).
\]
In the second limit the central summand uses \eqref{vieta}, the two
neighboring summands contribute $r^3$ each, and all remaining summands
vanish after division by $t^2$. Equation \eqref{moments} therefore
gives $h_2=(8f^2-r^2)/(12f)$.

Finally, the finite loci are complete lines. The unlabeled real
polygons form a circle, the quotient by $T$ of either real component
of $E$. Removing its unique infinite-vertex orbit leaves an interval.
For $t\to\pm\infty$, the original centroids satisfy
\[
 (C_0(P))_y=t/N+O(1),\qquad (C_2(P))_y=t/3+O(1).
\]
The second expansion follows from the same dominant triangle used
in \eqref{area-infty}. Equations \eqref{moments}--\eqref{vertex}
give the identical vertical coordinates for $Q$. They are continuous
on the finite interval and have opposite infinite limits at its
ends. Every real vertical coordinate is attained, completing the proof.

\section*{Verification and four proposed follow-up tasks}
The verifier checks the tangency and Vieta formulas, exact polygon
area/moment identities, pole cancellation, and numerical closed
families. The compact-curve argument is the proof for arbitrary $N$;
finite calculations are supplementary checks. The source's
perpendicular/directrix typo in its separate pentagon conjecture is
not used, and the perimeter-centroid degree is not asserted here.
\begin{enumerate}
 \item \textbf{Task 2.} Express $h_0$ in terms of $N$, $r/f$,
 and explicit elliptic-function invariants of the Poncelet map.
 \item \textbf{Task 3.} Determine the directions and offsets of the
 two polar centroid lines when the circular caustic center is displaced
 off the parabola's axis.
 \item \textbf{Task 4.} Determine the algebraic degree of the polar
 perimeter-centroid locus for every $N$, including the pentagon case.
 \item \textbf{Task 5.} Classify inner conics of a parabola for which
 the vertex and area centroids of the polar polygons trace parallel
 lines for every available finite Poncelet period.
\end{enumerate}
\begin{thebibliography}{9}
\bibitem{BGR} F. Bellio, R. Garcia, D. Reznik,
\emph{Parabola-Inscribed Poncelet Polygons Derived from the Bicentric Family},
KoG \textbf{26} (2022), 3--15,
\href{https://doi.org/10.31896/k.26.1}{doi:10.31896/k.26.1};
\href{https://arxiv.org/abs/2111.00979}{arXiv:2111.00979v3}.
\bibitem{ST} R. Schwartz, S. Tabachnikov,
\emph{Centers of mass of Poncelet polygons, 200 years after},
Math. Intelligencer \textbf{38}(2) (2016), 29--34,
\href{https://doi.org/10.1007/s00283-016-9622-9}{doi:10.1007/s00283-016-9622-9};
\href{https://arxiv.org/abs/1607.04766}{arXiv:1607.04766}.
\bibitem{CC} A. C. Chavez-Caliz,
\emph{More about areas and centers of Poncelet polygons},
Arnold Math. J. \textbf{7} (2021), 91--106,
\href{https://doi.org/10.1007/s40598-020-00154-8}{doi:10.1007/s40598-020-00154-8};
\href{https://arxiv.org/abs/2004.05404}{arXiv:2004.05404}.
\end{thebibliography}
\end{document}
